\section{Strong normalization for explicit polymorphism}
\label{sec:p07-system-f-normalization}

The normalization claim for System~F concerns the combined compatible
relation on Sheet~\ref{sheet:system-f}. Both term-beta and type-beta
contractions may occur at any position, including under either kind of
abstraction. The proof below applies to open typing judgments, permits
impredicative type instances, and retains the explicitly annotated term
syntax. It follows the candidate construction in
\cite[Chapter~14, with Chapters~6 and 11.1--11.2]{GirardLafontTaylor2003Web}.
The auxiliary candidates here contain raw Church-style terms; typing is
recovered through the fundamental lemma. This is a stated adaptation of
the source's candidates containing terms of a specified type.
The earlier simply typed construction in
Section~\ref{sec:p07-stlc-metatheory} supplies the arrow-only comparison.

\subsection{Syntax and substitution for the combined relation}
\label{sec:p07-system-f-sn-syntax}

\begin{hbdefinition}{def:p07-system-f-sn-raw}{Raw terms and normalization height}
Let \(\mathcal T_{\mathsf F}\) be the finite types generated by type
variables, \(\alpha\mathbin{\to_{\rm ty}}\beta\), and
\(\forall a.\alpha\). Let \(\mathcal M_{\mathsf F}\) be all finite
raw terms generated by
\[
 x,\quad MN,\quad\lambda x{:}\alpha.M,\quad
 \Lambda a.M,\quad M[\alpha],
\]
without requiring a typing derivation. Types and terms are taken modulo
consistent renaming of bound variables. A raw term is \emph{neutral} when
its outer constructor is a variable, term application or type application.
An application whose function is an abstraction is therefore neutral even
when it is a redex.

Write \(M\to_{\mathsf F}N\) for one contraction from the two beta rules
on the sheet, closed under all term constructors. Types have no reduction
rule: a type argument or annotation is replaced by type substitution but
is not itself reduced. Write \(\to_{\mathsf F}^*\) for finitely many
steps, including zero. Let \(\mathsf{SN}_{\mathsf F}\) be the set of
terms from which no infinite combined reduction starts. For a member
\(M\) of this set, \(h_{\mathsf F}(M)\) is the greatest possible
length of a reduction starting at \(M\).
A normal term has no immediate reduct. A maximal reduction sequence is
infinite or is finite and cannot be extended by another contraction.
In typing judgments below, \(\Delta\) abbreviates the sheet's
type-variable context \(\Delta_{\rm ty}\), and \(\Gamma\)
abbreviates its term-variable context \(\Theta_{\rm ty}\).
\index{System F!strong normalization}\index{reducibility candidate}
\end{hbdefinition}

\begin{hbdefinition}{def:p07-system-f-sn-substitutions}{The two simultaneous substitutions}
A finite type substitution \(\tau\) replaces free type variables by
members of \(\mathcal T_{\mathsf F}\), acting on all annotations and
type arguments in a term. A finite term substitution \(\sigma\) replaces
free term variables by raw terms. Both act once on the original syntax;
inserted images are not recursively substituted into. Bound variables of
both kinds are renamed when necessary to avoid capture in inserted images.
The notation \((M\tau)\sigma\) applies the type substitution first.
In particular, \(M[N/x]\) is term substitution and \(M[\beta/a]\)
is type substitution; \(M[\beta]\) without a slash is a type-application
constructor.

The finite sets \(\operatorname{FV}_{\rm tm}(M)\) and
\(\operatorname{FV}_{\rm ty}(M)\) record free term and type variables,
with free type variables counted in annotations and type arguments as well
as in term bodies. For a type \(\alpha\), \(\operatorname{FV}(\alpha)\)
denotes its free type variables. Substitution identities below are equalities modulo
alpha-renaming. Each time an identity crosses a binder, its name is
chosen fresh for every substitution image involved.
\end{hbdefinition}

\begin{hblemma}{thm:p07-system-f-sn-substitution-reduction}
 {proved=annex:p7.system-f.substitution-reduction}{Substitution preserves combined contractions}
For the raw syntax just defined,
\[
 M\to_{\mathsf F}M' \Longrightarrow\quad
 M\tau\to_{\mathsf F}M'\tau\quad\hbox{and}\quad
 M\sigma\to_{\mathsf F}M'\sigma.
\]
Also,
\[
 N\to_{\mathsf F}N' \Longrightarrow\quad
 M[N/x]\to_{\mathsf F}^*M[N'/x].
\]
The first two conclusions are genuine one-step contractions; the last
may use zero steps if there is no free occurrence of \(x\).
Every strongly normalizing term has finite height, and a one-step reduct
has strictly smaller height.
\end{hblemma}

\begin{hblemma}{thm:p07-system-f-sn-typing-substitution}
 {proved=annex:p7.system-f.typing-substitution}{Typing substitutions and preservation}
The following statements use exactly the contexts and rules of
Sheet~\ref{sheet:system-f}.
\begin{enumerate}
\item If \(\Delta;\Gamma\derives_{\mathsf F}M:\alpha\), and every
      image \(\tau(a)\), \(a\in\Delta\), is a type over a type context
      \(\Delta'\), then
      \(\Delta';\Gamma\tau\derives_{\mathsf F}M\tau:\alpha\tau\).
\item If \(\Delta;\Gamma\derives_{\mathsf F}M:\alpha\), and a common
      target context \(\Gamma'\) types every \(\sigma(x)\) at the
      type assigned to \(x\) in \(\Gamma\), then
      \(\Delta;\Gamma'\derives_{\mathsf F}M\sigma:\alpha\).
\item If \(\Delta;\Gamma\derives_{\mathsf F}M:\alpha\) and
      \(M\to_{\mathsf F}N\), then
      \(\Delta;\Gamma\derives_{\mathsf F}N:\alpha\).
\end{enumerate}
Here \(\Gamma\tau\) substitutes in each declared type. Weakening and
renaming use fresh variables and preserve the well-formedness conditions.
\end{hblemma}

\subsection{Candidates and impredicative interpretation}
\label{sec:p07-system-f-sn-candidates}

\begin{hbdefinition}{def:p07-system-f-sn-candidate}{Reducibility candidates for the raw combined calculus}
A candidate is a subset \(R\subseteq\mathcal M_{\mathsf F}\)
satisfying all three conditions:
\begin{enumerate}
\item \(R\subseteq\mathsf{SN}_{\mathsf F}\).
\item If \(M\in R\) and \(M\to_{\mathsf F}M'\), then \(M'\in R\).
\item If \(M\) is neutral and every immediate reduct of \(M\) is in
      \(R\), then \(M\in R\).
\end{enumerate}
Write \(\mathfrak C_{\mathsf F}\) for the set of all candidates.
It is a collection of subsets of one fixed set of raw terms, defined
before any interpretation of types. Clause~3 includes neutral normal
terms vacuously; in particular every variable belongs to every candidate.
Clause~2 iterates along finite reductions.
\end{hbdefinition}

\begin{hblemma}{thm:p07-system-f-sn-candidate-constructors}
 {proved=annex:p7.system-f.candidate-constructors}{Candidate formation}
The set \(\mathsf{SN}_{\mathsf F}\) is a candidate. If \(R,S\) are
candidates, so is
\[
 R\Rightarrow S=\{M:\text{for every }N\in R,\ MN\in S\}.
\]
For every family of candidates \(H_{\beta,C}\), indexed by all
\(\beta\in\mathcal T_{\mathsf F}\) and \(C\in\mathfrak C_{\mathsf F}\),
the set
\[
 \mathcal U(H)=\{M:\text{for every }\beta,C,\ M[\beta]\in H_{\beta,C}\}
\]
is a candidate. None of these definitions restricts type instances to
smaller, monomorphic or closed types.
\end{hblemma}

\begin{hbdefinition}{def:p07-system-f-sn-interpretation}{Interpretation with candidate parameters}
A candidate valuation \(\rho\) assigns a candidate to each type variable.
Define \(\mathcal R_\rho(\alpha)\) by structural recursion on the
finite type \(\alpha\):
\[
\begin{aligned}
 \mathcal R_\rho(a)&=\rho(a),\\
 \mathcal R_\rho(\alpha\mathbin{\to_{\rm ty}}\beta)
   &=\mathcal R_\rho(\alpha)\Rightarrow\mathcal R_\rho(\beta).
\end{aligned}
\]
For the universal type the defining condition is
\[
\begin{aligned}
 \mathcal R_\rho(\forall a.\alpha)=\{M:\quad&
       \forall\beta\in\mathcal T_{\mathsf F}\,
       \forall C\in\mathfrak C_{\mathsf F},\\
       &M[\beta]\in\mathcal R_{\rho[a\mapsto C]}(\alpha)\}.
\end{aligned}
\]
The valuation \(\rho[a\mapsto C]\) agrees with \(\rho\) except at
\(a\). The recursion at a universal type interprets the proper subtype
\(\alpha\) with a changed parameter. It does not recursively interpret
every syntactic instance \(\alpha[\beta/a]\). That difference removes
the circularity of the naive impredicative definition.
This is a set-theoretic interpretation involving sets of raw terms and
quantification over all candidates. It is not a claim that candidate
membership is decidable or that the construction is formalizable in a
specified weak arithmetic.
\end{hbdefinition}

\begin{hblemma}{thm:p07-system-f-sn-interpretation}
 {proved=annex:p7.system-f.interpretation}{Interpretations are candidates and respect type substitution}
For every valuation \(\rho\) and type \(\alpha\),
\(\mathcal R_\rho(\alpha)\) is a candidate. It depends on \(\rho\)
only at free variables of \(\alpha\), and is invariant under renaming
bound type variables. For every type \(\beta\),
\[
 \mathcal R_\rho(\alpha[\beta/a])
   =\mathcal R_{\rho[a\mapsto\mathcal R_\rho(\beta)]}(\alpha).
\]
The candidate \(\mathcal R_\rho(\beta)\) is available even when
\(\beta\) contains universal quantifiers, including a type as complicated
as the universal type being instantiated.
\end{hblemma}

\subsection{Introduction lemmas and the fundamental lemma}
\label{sec:p07-system-f-sn-fundamental}

\begin{hblemma}{thm:p07-system-f-sn-beta-expansion}
 {proved=annex:p7.system-f.beta-expansion}{Controlled expansion of the two beta redexes}
For a candidate \(R\), raw terms \(P,N\), and types \(\alpha,\beta\):
\begin{enumerate}
\item If \(P,N\in\mathsf{SN}_{\mathsf F}\) and \(P[N/x]\in R\),
      then \((\lambda x{:}\alpha.P)N\in R\).
\item If \(P\in\mathsf{SN}_{\mathsf F}\) and \(P[\beta/a]\in R\),
      then \((\Lambda a.P)[\beta]\in R\).
\end{enumerate}
These are guarded expansion statements. Arbitrary beta expansion need
not preserve strong normalization.
\end{hblemma}

\begin{hblemma}{thm:p07-system-f-sn-introduction}
 {proved=annex:p7.system-f.introduction}{Arrow and universal introduction in the interpretation}
If \(P[N/x]\in S\) for every \(N\in R\), where \(R,S\) are
candidates, then \(\lambda x{:}\alpha.P\in R\Rightarrow S\).
If, for every type \(\beta\) and candidate \(C\),
\[
 P[\beta/a]\in\mathcal R_{\rho[a\mapsto C]}(\alpha),
\]
then \(\Lambda a.P\in\mathcal R_\rho(\forall a.\alpha)\).
\end{hblemma}

\begin{hbtheorem}{thm:p07-system-f-sn-fundamental}
 {proved=annex:p7.system-f.fundamental}{Fundamental lemma with both substitutions}
Suppose \(\Delta;\Gamma\derives_{\mathsf F}M:\alpha\) on the
System~F sheet. Let \(\tau\) assign arbitrary finite System~F types
to the variables of \(\Delta\), and let \(\rho\) be any candidate
valuation. Let \(\sigma\) assign raw terms to the term variables of
\(\Gamma\), such that
\[
 x:\gamma\in\Gamma\quad\Longrightarrow\quad
              \sigma(x)\in\mathcal R_\rho(\gamma).
\]
Then
\[
             (M\tau)\sigma\in\mathcal R_\rho(\alpha).
\]
The images of \(\sigma\) need not be closed or typable: membership
in the indicated candidates is the exact hypothesis used by this
auxiliary lemma. The type substitution acts before the term substitution,
and the candidate valuation interprets the original types in the judgment.
\end{hbtheorem}

Taking the identity type substitution, the valuation assigning
\(\mathsf{SN}_{\mathsf F}\) to each type variable, and the identity
term substitution proves
Theorem~\ref{thm:p07-system-f-strong-normalization}, with its full open-term
scope. Its complete closing argument is
Annex~\ref{annex:p7.system-f.strong-normalization}.
The proof treats every possible contraction, so every maximal reduction
sequence is finite and ends at a normal term. The typing-preservation
lemma keeps its type unchanged. Confluence is the separate result
\ref{thm:p07-system-f-confluence}; normalization alone does not establish
uniqueness of the normal form.

\begin{hbexample}{def:p07-system-f-sn-example}{An impredicative instance and an open reduction}
Put \(I=\Lambda a.\lambda x{:}a.x\) and
\(\iota=\forall a.(a\mathbin{\to_{\rm ty}}a)\).
The sheet types \(I:\iota\), \(I[\iota]:\iota\mathbin{\to_{\rm ty}}\iota\),
and \(I[\iota]I:\iota\). Both kinds of contraction occur:
\[
 I[\iota]I\to_{\mathsf F}(\lambda x{:}\iota.x)I
                 \to_{\mathsf F}I.
\]
For a free type variable \(b\) and free term variable \(z:b\),
the term \(I[b]((\lambda y{:}b.y)z)\) is typable at \(b\).
One may reduce its argument first or contract the type application first;
in either case the displayed redexes eventually give \(z\).
The theorem allows these open variables throughout the reduction.
\end{hbexample}

\begin{hbwarning}{def:p07-system-f-sn-side-condition}{Why the hypotheses cannot be dropped}
From \(a;x:a\derives_{\mathsf F}x:a\), one cannot infer a judgment
typing \(\Lambda a.x\) at \(\forall a.a\) with that same free
term declaration: \(a\) occurs in its declared type. In the fundamental
lemma, an image known to belong to \(\rho(a)\) cannot simply be
assumed to belong to every candidate \(C\). The eigenvariable condition
is what keeps all context interpretations unchanged during the universal
introduction case.

The strong-normalization hypotheses in beta expansion are also necessary.
For a fixed type variable \(a\), the raw term
\(\Omega=(\lambda x{:}a.xx)(\lambda x{:}a.xx)\) contracts to itself.
The raw application \((\lambda u{:}a.z)\Omega\) has normal root
contractum \(z\), but it admits an infinite sequence reducing only its
argument. These raw terms are not offered as well-typed System~F terms.
They show why candidate expansion must control the unreduced body and
argument, rather than infer termination from one terminating contraction.
\end{hbwarning}
